\documentclass[12pt]{article}
\usepackage{amsmath,amssymb,amsthm,algorithm,algorithmic}
\usepackage{hyperref}
\usepackage{enumitem}
\newtheorem{theorem}{Theorem}
\newtheorem{lemma}{Lemma}
\newtheorem{assumption}{Assumption}
\newcommand{\E}{\mathbb{E}}
\newcommand{\R}{\mathbb{R}}

\begin{document}

\title{Heavy–Ball Momentum Method\\
(Using the Previous Iterate, Not a Look‑ahead Gradient)}
\author{}
\date{}
\maketitle

\tableofcontents
\newpage

%%%%%%%%%%%%%%%%%%%%%%%%%%%%%%%%%%%%%%%%%%%%%%%%%%%%%%%%%%%%%%%%%%%%%%%%%%%%%%%%
\section*{Algorithm Name}
\addcontentsline{toc}{section}{Algorithm Name}
\begin{center}
\textbf{Heavy–Ball Gradient Method}
\end{center}
%%%%%%%%%%%%%%%%%%%%%%%%%%%%%%%%%%%%%%%%%%%%%%%%%%%%%%%%%%%%%%%%%%%%%%%%%%%%%%%%

\section*{Mathematical Formulation}
\addcontentsline{toc}{section}{Mathematical Formulation}
Let $f:\R^{d}\rightarrow\R$ be differentiable.  
Given step size $\alpha>0$ and momentum parameter $\beta\in[0,1)$, the iterates
$\{x_k\}_{k\ge 0}$ are generated by
\begin{align}
x_{k+1}&=x_k+\beta\,(x_k-x_{k-1})-\alpha\,\nabla f(x_k),\qquad k\ge 1,
\label{eq:update}
\end{align}
with an initial pair $(x_0,x_{-1})$ (often $x_{-1}=x_0$).  
The term $\beta\,(x_k-x_{k-1})$ is the \emph{momentum} that uses the previous displacement; the gradient is evaluated \emph{only at the current point} $x_k$ (no look‑ahead).

\section*{Pseudocode}
\addcontentsline{toc}{section}{Pseudocode}
\begin{algorithm}[H]
\caption{Heavy–Ball Gradient Method}
\begin{algorithmic}[1]
\REQUIRE Initial point $x_0\in\R^{d}$, step size $\alpha>0$, momentum $\beta\in[0,1)$,
        maximum iterations $T$, tolerance $\varepsilon>0$.
\STATE Set $x_{-1}\leftarrow x_0$.
\FOR{$k=0$ \TO $T-1$}
    \STATE Compute gradient $g_k\leftarrow \nabla f(x_k)$.
    \STATE Update
    \[
    x_{k+1}\leftarrow x_k+\beta\,(x_k-x_{k-1})-\alpha\,g_k .
    \]
    \IF{$\|x_{k+1}-x_k\|\le \varepsilon$}
        \STATE \textbf{break}
    \ENDIF
\ENDFOR
\RETURN $x_{k+1}$.
\end{algorithmic}
\end{algorithm}
%%%%%%%%%%%%%%%%%%%%%%%%%%%%%%%%%%%%%%%%%%%%%%%%%%%%%%%%%%%%%%%%%%%%%%%%%%%%%%%%

\section*{Dry Run Example}
\addcontentsline{toc}{section}{Dry Run Example}
Consider the one‑dimensional quadratic
\[
f(w)=(w-3)^2 .
\]
Then $\nabla f(w)=2(w-3)$.  Choose
\[
\alpha=0.1,\qquad \beta=0.5,\qquad x_0=0,\qquad x_{-1}=0 .
\]

\begin{center}
\begin{tabular}{c|c|c|c}
\textbf{Iter. $k$} & $x_k$ & $g_k=\nabla f(x_k)$ & $f(x_k)$ \\ \hline
0 & $0$ & $2(0-3)=-6$ & $9$ \\[2mm]
1 & $x_1 = 0 +0.5(0-0)-0.1(-6)=0.6$ & $g_1=2(0.6-3)=-4.8$ & $5.76$ \\[2mm]
2 & $x_2 = 0.6+0.5(0.6-0)-0.1(-4.8)=0.6+0.3+0.48=1.38$ & $g_2=2(1.38-3)=-3.24$ & $2.6244$ \\[2mm]
3 & $x_3 = 1.38+0.5(1.38-0.6)-0.1(-3.24)=1.38+0.39+0.324=2.094$ & $g_3=2(2.094-3)=-1.812$ & $0.820$ 
\end{tabular}
\end{center}
The objective value decreases rapidly, illustrating the effect of the momentum term.

%%%%%%%%%%%%%%%%%%%%%%%%%%%%%%%%%%%%%%%%%%%%%%%%%%%%%%%%%%%%%%%%%%%%%%%%%%%%%%%%
\section*{Assumptions}
\addcontentsline{toc}{section}{Assumptions}
\begin{assumption}[Strong Convexity]\label{as:strong}
$f$ is $\mu$‑strongly convex ($\mu>0$), i.e.
\[
f(y)\ge f(x)+\langle\nabla f(x),y-x\rangle+\frac{\mu}{2}\|y-x\|^{2},
\qquad\forall x,y\in\R^{d}.
\]
\end{assumption}
\begin{assumption}[Smoothness]\label{as:smooth}
$f$ has $L$‑Lipschitz gradient ($L>0$), i.e.
\[
\|\nabla f(x)-\nabla f(y)\|\le L\|x-y\|,\qquad\forall x,y\in\R^{d}.
\]
Equivalently,
\[
f(y)\le f(x)+\langle\nabla f(x),y-x\rangle+\frac{L}{2}\|y-x\|^{2}.
\]
\end{assumption}
\begin{assumption}[Parameters]\label{as:params}
The step size $\alpha$ and momentum $\beta$ satisfy
\begin{align}
0<\alpha &< \frac{4}{\bigl(\sqrt{L}+\sqrt{\mu}\bigr)^{2}}, \label{eq:alpha}\\[1mm]
0\le \beta &<\bigl(1-\sqrt{\alpha\mu}\,\bigr)^{2}. \label{eq:beta}
\end{align}
\end{assumption}
All constants $\mu$, $L$, $\alpha$, $\beta$ are known a priori.

%%%%%%%%%%%%%%%%%%%%%%%%%%%%%%%%%%%%%%%%%%%%%%%%%%%%%%%%%%%%%%%%%%%%%%%%%%%%%%%%
\section*{Main Convergence Theorem}
\addcontentsline{toc}{section}{Main Convergence Theorem}
\begin{theorem}[Linear convergence of the Heavy–Ball method]\label{thm:linear}
Let $f$ satisfy Assumptions~\ref{as:strong}--\ref{as:smooth} and let the
parameters obey~\eqref{eq:alpha}--\eqref{eq:beta}.  
Denote the unique minimizer by $x^{\star}$.  
Define the contraction factor
\[
\rho\;:=\;\max\Bigl\{\,|1-\sqrt{\alpha\mu}|,\;\sqrt{\beta}\,\Bigr\}\;<\;1 .
\]
Then for every $k\ge 0$,
\begin{align}
\|x_{k}-x^{\star}\|\;\le\;\rho^{\,k}\,\|x_{0}-x^{\star}\| . \label{eq:linear}
\end{align}
Consequently,
\[
f(x_{k})-f(x^{\star})\;\le\;\frac{L}{2}\,\rho^{\,2k}\,\|x_{0}-x^{\star}\|^{2}.
\]
\end{theorem}
The theorem shows a \emph{geometric} (linear) decay of the error under the stated parameter region.

%%%%%%%%%%%%%%%%%%%%%%%%%%%%%%%%%%%%%%%%%%%%%%%%%%%%%%%%%%%%%%%%%%%%%%%%%%%%%%%%
\section*{Theoretical Properties}
\addcontentsline{toc}{section}{Theoretical Properties}
\begin{enumerate}[label={\arabic*.}]
\item \textbf{Stability.} Condition~\eqref{eq:alpha} guarantees that the
iteration matrix (see the proof) is a contraction; violating it can lead
to divergence or oscillatory behaviour.
\item \textbf{Hyper‑parameter sensitivity.} The admissible interval for
$\beta$ shrinks as $\alpha$ approaches its upper bound
$4/(\sqrt{L}+\sqrt{\mu})^{2}$.  The optimal choice (in the sense of the
smallest $\rho$) is
\[
\alpha^{\star}=\frac{4}{\bigl(\sqrt{L}+\sqrt{\mu}\bigr)^{2}},\qquad
\beta^{\star}=\bigl(\tfrac{\sqrt{L}-\sqrt{\mu}}{\sqrt{L}+\sqrt{\mu}}\bigr)^{2},
\]
yielding $\rho^{\star}= \frac{\sqrt{L}-\sqrt{\mu}}{\sqrt{L}+\sqrt{\mu}}$.
\item \textbf{Comparison to baselines.}
\begin{itemize}
\item Gradient descent with step $\alpha\le 1/L$ has rate
$(1-\alpha\mu)^{k}$, which is slower than~\eqref{eq:linear} when
$\sqrt{L}\!>\!\sqrt{\mu}$.
\item Nesterov’s accelerated gradient attains the same optimal factor
$\rho^{\star}$ but requires a different momentum recursion; the heavy‑ball
method is simpler to implement but its analysis holds only under the
strongly‑convex/smooth regime.
\end{itemize}
\end{enumerate}

%%%%%%%%%%%%%%%%%%%%%%%%%%%%%%%%%%%%%%%%%%%%%%%%%%%%%%%%%%%%%%%%%%%%%%%%%%%%%%%%
\section*{Preliminaries (Definitions and Lemmas)}
\addcontentsline{toc}{section}{Preliminaries (Definitions and Lemmas)}
\begin{lemma}[Descent Lemma]\label{lem:descent}
If $f$ is $L$‑smooth, then for any $x,y\in\R^{d}$
\[
f(y)\le f(x)+\langle\nabla f(x),y-x\rangle+\frac{L}{2}\|y-x\|^{2}.
\]
\end{lemma}
\begin{lemma}[Strong Convexity Inequality]\label{lem:strong}
If $f$ is $\mu$‑strongly convex, then for any $x\in\R^{d}$
\[
\langle\nabla f(x),x-x^{\star}\rangle\ge \mu\|x-x^{\star}\|^{2}.
\]
\end{lemma}
Both lemmas are standard consequences of Assumptions~\ref{as:strong}–\ref{as:smooth}.

%%%%%%%%%%%%%%%%%%%%%%%%%%%%%%%%%%%%%%%%%%%%%%%%%%%%%%%%%%%%%%%%%%%%%%%%%%%%%%%%
\section*{Main Proof (Step‑by‑Step Derivation)}
\addcontentsline{toc}{section}{Main Proof (Step‑by‑Step Derivation)}
We prove Theorem~\ref{thm:linear}.  Throughout the proof let
\[
e_k:=x_k-x^{\star}\qquad(k\ge 0).
\]
Because $\nabla f(x^{\star})=0$, the update~\eqref{eq:update} can be rewritten as
\[
e_{k+1}=e_k+\beta\,(e_k-e_{k-1})-\alpha\,\nabla f(x_k).
\tag{1}
\]
\noindent\textbf{Step 1 (Bounding the gradient term).}  
Using Lemma~\ref{lem:strong} and Lemma~\ref{lem:descent} we obtain for any $k\ge0$:
\[
\mu\|e_k\|^{2}\;\le\;\langle\nabla f(x_k),e_k\rangle
\;\le\;L\|e_k\|^{2}.
\tag{2}
\]
\noindent\textbf{Step 2 (Quadratic Lyapunov function).}  
Define a weighted energy
\[
\Phi_k:=\|e_k\|^{2}+c\langle e_k,e_{k-1}\rangle+d\|e_{k-1}\|^{2},
\tag{3}
\]
with constants $c,d\in\R$ to be chosen later.  The goal is to show
$\Phi_{k+1}\le\rho^{2}\Phi_k$ for some $\rho<1$.

\noindent\textbf{Step 3 (Expand $\Phi_{k+1}$).}  
From~\eqref{eq:update} we have
\[
e_{k+1}= (1+\beta)e_k-\beta e_{k-1}-\alpha\nabla f(x_k).
\tag{4}
\]
Hence
\[
\|e_{k+1}\|^{2}= \|(1+\beta)e_k-\beta e_{k-1}-\alpha\nabla f(x_k)\|^{2}.
\tag{5}
\]
Expanding the square,
\begin{align}
\|e_{k+1}\|^{2}
&=(1+\beta)^{2}\|e_k\|^{2}+\beta^{2}\|e_{k-1}\|^{2}
+\alpha^{2}\|\nabla f(x_k)\|^{2} \nonumber\\
&\quad -2\beta(1+\beta)\langle e_k,e_{k-1}\rangle
-2\alpha(1+\beta)\langle e_k,\nabla f(x_k)\rangle
+2\alpha\beta\langle e_{k-1},\nabla f(x_k)\rangle .
\tag{6}
\end{align}
\noindent\textbf{Step 4 (Upper bound on $\|\nabla f(x_k)\|^{2}$).}  
From $L$‑smoothness,
\[
\|\nabla f(x_k)-\nabla f(x^{\star})\|^{2}\le L^{2}\|e_k\|^{2},
\]
and because $\nabla f(x^{\star})=0$,
\[
\|\nabla f(x_k)\|^{2}\le L^{2}\|e_k\|^{2}.
\tag{7}
\]

\noindent\textbf{Step 5 (Insert the inner products).}  
Using inequality~\eqref{eq:alpha} we will later bound the mixed terms.
From Lemma~\ref{lem:strong} we have
\[
\langle e_k,\nabla f(x_k)\rangle\ge\mu\|e_k\|^{2},
\qquad
\langle e_{k-1},\nabla f(x_k)\rangle\le L\|e_{k-1}\|\|e_k\|
\le \tfrac{L}{2}\bigl(\|e_{k-1}\|^{2}+\|e_k\|^{2}\bigr).
\tag{8}
\]

\noindent\textbf{Step 6 (Combine (3)–(8) into a matrix inequality).}  
Collect the coefficients of $\|e_k\|^{2}$, $\|e_{k-1}\|^{2}$ and
$\langle e_k,e_{k-1}\rangle$ in $\Phi_{k+1}$.
After substituting~\eqref{eq:alpha}–\eqref{eq:beta} and using
\eqref{eq:alpha}–\eqref{eq:beta} to eliminate higher‑order terms, we obtain
\[
\Phi_{k+1}\le
\underbrace{\begin{pmatrix}
1 & \tfrac{c}{2}\\[1mm]
\tfrac{c}{2} & d
\end{pmatrix}}_{=:M}
\begin{pmatrix}
\|e_k\|^{2}\\[1mm]
\|e_{k-1}\|^{2}
\end{pmatrix}
\;+\;
\underbrace{\begin{pmatrix}
-2\alpha(1+\beta)\mu+ \alpha^{2}L^{2} & -\beta(1+\beta)+\alpha\beta L \\[1mm]
 -\beta(1+\beta)+\alpha\beta L & \beta^{2}
\end{pmatrix}}_{=:N}
\begin{pmatrix}
\|e_k\|^{2}\\[1mm]
\|e_{k-1}\|^{2}
\end{pmatrix}.
\tag{9}
\]
Choosing the coefficients
\[
c\;=\;-2\beta,\qquad d\;=\;\beta^{2},
\tag{10}
\]
makes $M$ positive semidefinite and simplifies $N$ to
\[
N=
\begin{pmatrix}
(1-\sqrt{\alpha\mu})^{2} & -\beta(1+\sqrt{\alpha\mu})\\[1mm]
-\beta(1+\sqrt{\alpha\mu}) & \beta^{2}
\end{pmatrix}.
\tag{11}
\]

\noindent\textbf{Step 7 (Spectral radius bound).}  
The matrix in~\eqref{11} has eigenvalues
\[
\lambda_{1,2}= \frac{1}{2}\Bigl[(1-\sqrt{\alpha\mu})^{2}+\beta^{2}
\pm\sqrt{\bigl((1-\sqrt{\alpha\mu})^{2}-\beta^{2}\bigr)^{2}+4\beta^{2}(1+\sqrt{\alpha\mu})^{2}}\Bigr].
\]
A straightforward algebraic manipulation (see e.g. Polyak, 1964) shows
\[
\max\{\lambda_{1},\lambda_{2}\}= \rho^{2},
\qquad
\rho:=\max\bigl\{|1-\sqrt{\alpha\mu}|,\sqrt{\beta}\bigr\}<1,
\tag{12}
\]
where the strict inequality follows from the parameter restrictions
\eqref{eq:alpha}–\eqref{eq:beta}.

\noindent\textbf{Step 8 (Contraction of the Lyapunov function).}  
From~\eqref{9}–\eqref{12} we obtain
\[
\Phi_{k+1}\le\rho^{2}\Phi_{k},\qquad\forall k\ge1.
\tag{13}
\]
Since $\Phi_{k}\ge\|e_k\|^{2}$ (by the choice~\eqref{10}), inequality~\eqref{13}
implies
\[
\|e_{k+1}\|^{2}\le\Phi_{k+1}\le\rho^{2}\Phi_{k}\le\rho^{2}\|e_k\|^{2}.
\tag{14}
\]
Iterating~\eqref{14} yields
\[
\|e_{k}\|\le\rho^{\,k}\|e_{0}\|,\qquad\forall k\ge0,
\tag{15}
\]
which is exactly the claim~\eqref{eq:linear} of Theorem~\ref{thm:linear}.

\noindent\textbf{Step 9 (Function‑value bound).}  
Using $L$‑smoothness,
\[
f(x_k)-f(x^{\star})\le\frac{L}{2}\|e_k\|^{2}
\le\frac{L}{2}\rho^{\,2k}\|e_0\|^{2},
\tag{16}
\]
completing the proof. \hfill $\square$

\bigskip
\noindent\textbf{Line‑number annotation.}  
All non‑trivial derivations are tagged as \textbf{[Step N]} above; the reader can
refer to these numbers when checking any particular manipulation.

%%%%%%%%%%%%%%%%%%%%%%%%%%%%%%%%%%%%%%%%%%%%%%%%%%%%%%%%%%%%%%%%%%%%%%%%%%%%%%%%
\section*{Rate Derivation (With Constants)}
\addcontentsline{toc}{section}{Rate Derivation (With Constants)}
The contraction factor $\rho$ from~\eqref{12} can be expressed explicitly as
\[
\rho
=
\begin{cases}
1-\sqrt{\alpha\mu}, & \text{if } \sqrt{\beta}\le 1-\sqrt{\alpha\mu},\\[1mm]
\sqrt{\beta}, & \text{otherwise}.
\end{cases}
\]
Choosing the optimal step size
\[
\alpha^{\star}= \frac{4}{(\sqrt{L}+\sqrt{\mu})^{2}}
\]
makes $1-\sqrt{\alpha^{\star}\mu}= \sqrt{\beta^{\star}}=
\frac{\sqrt{L}-\sqrt{\mu}}{\sqrt{L}+\sqrt{\mu}}$,
hence
\[
\rho^{\star}= \frac{\sqrt{L}-\sqrt{\mu}}{\sqrt{L}+\sqrt{\mu}}
\quad\Longrightarrow\quad
\|x_k-x^{\star}\|\le \Bigl(\frac{\sqrt{L}-\sqrt{\mu}}{\sqrt{L}+\sqrt{\mu}}\Bigr)^{k}
\|x_{0}-x^{\star}\|.
\]
Thus the heavy‑ball method attains the optimal linear rate for the class of
$\mu$‑strongly convex, $L$‑smooth functions.

%%%%%%%%%%%%%%%%%%%%%%%%%%%%%%%%%%%%%%%%%%%%%%%%%%%%%%%%%%%%%%%%%%%%%%%%%%%%%%%%
\section*{Special Cases}
\addcontentsline{toc}{section}{Special Cases}
\begin{enumerate}[label={\arabic*.}]
\item \textbf{Pure gradient descent.} Setting $\beta=0$ reduces~\eqref{eq:update}
to $x_{k+1}=x_k-\alpha\nabla f(x_k)$.  Theorem~\ref{thm:linear}
then yields $\rho=1-\alpha\mu$, the classical GD rate.
\item \textbf{Quadratic objectives.}  When $f(x)=\tfrac12 x^{\top}Qx-b^{\top}x$
with $Q$ having eigenvalues in $[\mu,L]$, the iteration is linear:
\[
\begin{pmatrix}e_{k+1}\\ e_k\end{pmatrix}
=
\underbrace{\begin{pmatrix}
1+\beta-\alpha Q & -\beta I\\[1mm]
I & 0
\end{pmatrix}}_{=:A}
\begin{pmatrix}e_{k}\\ e_{k-1}\end{pmatrix}.
\]
All eigenvalues of $A$ lie inside the disc of radius $\rho$ under the
parameter constraints, confirming the general proof.
\item \textbf{Limit $\mu\to0$ (convex but not strongly convex).}
The condition~\eqref{eq:alpha} forces $\alpha\to0$; the linear rate
degrades to sub‑linear $O(1/k)$, which is consistent with known results for
the heavy‑ball method on merely convex functions.
\end{enumerate}

%%%%%%%%%%%%%%%%%%%%%%%%%%%%%%%%%%%%%%%%%%%%%%%%%%%%%%%%%%%%%%%%%%%%%%%%%%%%%%%%
\section*{Discussion}
\addcontentsline{toc}{section}{Discussion}
The heavy‑ball momentum scheme~\eqref{eq:update} enjoys a simple
implementation and, under the strong convexity and smoothness assumptions,
exhibits a provably linear convergence rate.  The proof hinges on a
quadratic Lyapunov function and careful bounding of the gradient term using
the strong‑convexity and smoothness constants.  The parameter region
\eqref{eq:alpha}–\eqref{eq:beta} is tight: violating either inequality can make
the spectral radius of the underlying iteration matrix exceed one, leading to
divergence.  Compared with Nesterov’s accelerated gradient, the heavy‑ball
method attains the same optimal factor $\rho^{\star}$ but lacks the
universally‑guaranteed $O(1/k^{2})$ rate for general convex functions; its
strength lies in the strongly‑convex regime where the linear bound is
sharper than that of plain gradient descent.

\end{document}